\clearpage
\section{感知机：从分错的样本修正线性边界}
\label{sec:linear-perceptron}
感知机（Perceptron）使用得分$s=\vct w^\top\vct x+b$区分两个类别，
标签约定为$y\in\{-1,+1\}$。边界为$s=0$；这里预测时$s\geq0$判$+1$。
它直接学习分类边界，不输出已经校准的概率。

\begin{figure}[H]
\centering\begin{tikzpicture}
\begin{axis}[width=120mm,height=49mm,xmin=-2.7,xmax=2.7,ymin=-2.5,ymax=2.7,
 axis lines=middle,xlabel={$x_1$},ylabel={$x_2$},xtick={-2,-1,1,2},ytick={-2,-1,1,2},
 legend style={font=\scriptsize,at={(1.02,1)},anchor=north west,draw=none}]
\addplot[only marks,mark=*,Teal,mark size=2.7pt] coordinates {(2,1) (-1,2)};
\addlegendentry{$y=+1$}
\addplot[only marks,mark=triangle*,Amber,mark size=3pt] coordinates {(1,-1) (-2,-1)};
\addlegendentry{$y=-1$}
\addplot[Amber,dashed,thick,domain=-1.8:.7] {-2*x-1};
\addlegendentry{$2x_1+x_2+1=0$}
\addplot[Teal,thick,domain=-2.7:2.7] {-.5*x};
\addlegendentry{$x_1+2x_2=0$}
\end{axis}
\end{tikzpicture}

\caption{四个可分点。虚线是第一次更新后的边界，实线是第二次更新后的边界。}
\end{figure}
\readfigure{依次读取$(2,1),+1$与$(1,-1),-1$。
从$\vct w=(0,0),b=0$开始，步长1，第一次更新为$(2,1),1$；
第二条得分为2，符号错误，更新后为$(1,2),0$。此时四点都严格分对。}

本页采用在线规则：当$y_i(\vct w^\top\vct x_i+b)\leq0$时更新
\[
 \vct w\leftarrow\vct w+\eta y_i\vct x_i,\qquad
 b\leftarrow b+\eta y_i,\quad\eta>0;
\]
否则保持。边界点即使按预测平票规则判对，也按此训练规则更新。
可把输入增广为$(\vct x,1)$，将偏置一并理解为参数。

\begin{pitfall}[有限步结论需要可分条件]
对有界、以正间隔线性可分的固定训练数据，经典感知机规则有有限次更新的收敛保证。
非线性可分或含冲突标签时可能不停更新，需要训练上限或改用合适目标。
收敛只说明训练集分开，不保证最大间隔、概率校准或未来表现。
\end{pitfall}

阶跃输出几乎处处导数为0，不能把上述更新理解为直接穿过硬阶跃的普通梯度。
可从感知机损失$\max(0,-ys)$的次梯度解释更新；在$ys=0$处需选定约定。
逻辑回归使用平滑的逻辑损失，SVM使用间隔与惩罚；三者不能只因边界线性就混为同一算法。


\noindent\textbf{继续阅读：}Shalev-Shwartz与Ben-David，\textit{Understanding Machine Learning}（2014），第9章。
